% Generated by tools/build_v5.py; edit the assembly rules, not this file.
\section{Model, notation, and main results}\label{sec:model}
Fix $a=1/2$. Let $\mu$ be any probability measure on $(0,a]$; in particular,
$\mu$ may have atoms, gaps, or a singular continuous part. A measurable
function $q:(0,a]\to[0,1]$ specifies a population by
\begin{equation}\label{eq:model}
 D\sim\mu,\qquad S\mid D=d\sim\Ber(q(d)),\qquad
 P=\tfrac12+(2S-1)D.
\end{equation}
The variables attached to distinct new coins are independent. The
population class is $\cG(\mu)=\{G_{\mu,q}:q\text{ measurable}\}$ and
$\theta(\mu,q)=\int q(d)\,\mu(\mathrm d d)$.
Individual magnitudes and signs are not observed. The upper algorithm does
not know $\mu$ either. The distribution-wise minimax benchmark may grant
knowledge of $\mu$, so attaining it simultaneously without that knowledge
is the stronger statement. The lower bound remains valid with still more
information. The endpoint $D=1/2$ is allowed and corresponds to deterministic
coins; $D=0$ is excluded.

An algorithm has an independent random seed, creates finitely many coins
before each observation, selects one of its coins, and receives its next
Bernoulli observation. It stops after $T$ observations and returns a real
number $\widehat\theta$. Infinite expected work is permitted in the model,
but algorithms satisfying a finite uniform budget necessarily stop almost
surely. Finite internal computation between observations is implicit.
We count observations, not arithmetic operations or independent random bits.
Every coin actually created by the upper algorithm receives at least one
observation. Charging a fixed additional amount per new outer draw therefore
changes its cost by at most a constant factor.

Define
\begin{align}
 R_\mu(B)&=\inf_{\cA:\,\sup_q\E_{\mu,q}T\le B}
       \sup_q\bigl(\E_{\mu,q}(\widehat\theta-\theta)^2\bigr)^{1/2},
       \label{eq:risk}\\
 r_\mu(B,\delta)&=\inf_{\cA:\,\sup_q\E_{\mu,q}T\le B}
   \inf\{r\ge0:\sup_q\Prob_{\mu,q}(|\widehat\theta-\theta|>r)\le\delta\}.
       \label{eq:confidence-risk}
\end{align}
These are expected-budget criteria, not deterministic maximum-work criteria.

For $h>0$, write
\begin{equation}\label{eq:A-I}
 \Am(h)=\E_\mu\min\{h^{-1},D^{-1}\},\qquad
 \om(h)=h\Am(h),\qquad \Iinv=\E_\mu[D^{-1}]\in[2,\infty].
\end{equation}
The modulus is increasing and concave, $\om(0+)=0$, and $\om(h)=1$ for
$h\ge a$, while $\om(h)\ge\min\{1,2h\}$. In particular, for $c\ge1$ and $0<b\le1$,
\begin{equation}\label{eq:scaling}
 \om(ch)\le c\om(h),\qquad \om(bh)\ge b\om(h).
\end{equation}

\begin{theorem}[Distribution-blind arbitrary-margin minimax law]\label{thm:main}
There are universal constants $0<c<C<\infty$ such that for every probability
measure $\mu$ on $(0,1/2]$, every $B\ge1$, and every $0<\delta\le1/4$,
\begin{align}
 c\om(B^{-1/2})&\le R_\mu(B)\le C\om(B^{-1/2}),\label{eq:main-rms}\\
 c\om\!\left(\sqrt{\frac{\log(1/\delta)}{B}}\right)
 &\le r_\mu(B,\delta)
 \le C\om\!\left(\sqrt{\frac{\log(1/\delta)}{B}}\right).
 \label{eq:main-conf}
\end{align}
More strongly, for each $B$ there is one algorithm $\cA_B$, depending only
on $B$, that simultaneously for all $\mu,q$ satisfies
\begin{equation}\label{eq:blind-guarantee}
 \E_{\mu,q}T\le B,\qquad
 \bigl(\E_{\mu,q}(\widehat\theta-\theta)^2\bigr)^{1/2}
 \le C\om(B^{-1/2}).
\end{equation}
There is likewise one algorithm depending only on $(B,\delta)$ for the
upper confidence bound. Neither algorithm observes $D$ or knows $\mu$.
In the lower bounds one may take $c=1/64$ for RMS and
$c=15/(256\sqrt2)$ for confidence.
The lower bounds hold even if every coin's magnitude is revealed for free,
or if the algorithm selects the magnitude of each fresh coin. In the latter
experiment $q$ is defined on all $(0,1/2]$ and the target is still
$\int q\,\mathrm d\mu$.
\end{theorem}

\paragraph{Formal scope.}
The RMS law in the ordinary hidden-magnitude experiment is proved by
\leanref{Tier6/Full.lean}{minimax\_lower\_bound} and
\leanref{BlindUpper/Full.lean}{distribution\_blind\_upper}, with
coefficients $1/64$ and $512$. The upper family is chosen before every
gap law and satisfies the budget in the same experiment. The
confidence statement and stronger
revealed- or selected-magnitude claims retain paper-proof status.

This is distribution-wise minimax optimality: the worst case is over all
sign functions at a fixed magnitude law. It is not instance optimality
for each particular pair $(\mu,q)$. The budget is expected work, not an
almost-sure cap. The modulus in the guarantee is evaluated at the true,
possibly unknown distribution; it is not an observable confidence radius.

\begin{theorem}[Exact unbiased feasibility and variance--work scale]
\label{thm:unbiased}
An estimator unbiased for every $G\in\cG(\mu)$ can have both
\[
 V=\sup_q\Var_{\mu,q}(Y)<\infty,\qquad
 W=\sup_q\E_{\mu,q}T<\infty
\]
if and only if $\Iinv<\infty$. When $\Iinv<\infty$,
\begin{equation}\label{eq:product}
 \frac1{16}\,\Iinv^2
 \ \le\ \inf_{Y\text{ unbiased}} VW
 \ \le\ 713\Iinv^2,
\end{equation}
where the infimum is over estimators with both uniform bounds finite.
A single estimator, using neither $\mu$ nor any margin parameters, attains
the upper bound simultaneously for every such magnitude law. It satisfies
$\sup_q\E Y^2\le23\Iinv$ and $\sup_q\E T\le31\Iinv$.
The lower coefficient $1/16$ is the exact universal constant: equality is
attained for every fixed-gap law $\mu=\delta_d$. Equality is not asserted
for arbitrary nondegenerate $\mu$.
\end{theorem}

\noindent\emph{Verification: Lean proves the full existence characterization,
the $1/16$ lower bound, and the 713 upper construction in the same
observable model; see \cref{sec:713,sec:verification}.}

\subsection{Unknown magnitude laws under an envelope}
Let $F_\mu(t)=\mu((0,t])$, extended to equal one above $a$, and define
\begin{equation}\label{eq:envelope-class}
 \cG_{\preceq\mu}=
 \{G\text{ on }[0,1]:\Prob_G(|P-1/2|\le t)\le F_\mu(t)
                    \text{ for every }t>0\}.
\end{equation}
Here the actual magnitude law is unknown on $(0,1/2]$.
The inequality means its magnitudes are stochastically no smaller than
those under $\mu$. It also rules out an atom at $P=1/2$.

\begin{corollary}[Stochastic-dominance envelope version]\label[corollary]{cor:envelope}
Both conclusions of \cref{thm:main} hold with the supremum taken over
$\cG_{\preceq\mu}$, with the same universal constants. The existence criterion
and variance--work product bounds of \cref{thm:unbiased} also hold for this
class. The finite-budget and unbiased algorithms need not know the envelope;
it is used only to express a uniform guarantee.
\end{corollary}

\paragraph{What distribution-blindness does not supply.}
The algorithm is specified without $\mu$ or an envelope. A numerical choice
of budget for a prescribed accuracy, however, generally requires a bound on
the unknown modulus. Neither the adaptive risk guarantee nor the unbiased
estimator automatically supplies a data-computable confidence radius.

\begin{corollary}[Critical nested-simulation obstruction]\label[corollary]{cor:nested}
Let $U\sim\mathrm{Unif}[0,1]$, $X\in\{-1,1\}$, and
$D(U)=|\E[X\mid U]|/2$. Consider all measurable conditional models with
$D(U)\sim\mathrm{Unif}(0,1/4]$, with no cross-state regularity.
For estimating $\Prob(\E[X\mid U]>0)$, the minimax expected observation
cost at RMS error $\varepsilon$, or at error $\varepsilon$ with a fixed
failure probability $\delta_0=1/16$, is
\[
 \asymp\varepsilon^{-2}\log^2(1/\varepsilon),\qquad \varepsilon\downarrow0.
\]
The lower bound allows arbitrary selected-state queries and revisiting.
Every model satisfies Assumptions 2.1 and 2.4 of Giles--Haji-Ali
\cite{gileshajiali}, with uniform constants.
\end{corollary}
The continuous-magnitude construction and proof are given in
\cref{prop:continuous}. This is a worst-case sample-oracle statement on a
subclass satisfying those assumptions, not a claim about every structured
simulator or every possible source of information.

\section{A history-uniform adaptive information bound}\label{sec:lower}
For Bernoulli probabilities $u,v\in(0,1)$,
\begin{equation}\label{eq:ber-kl}
 \kl(u,v)\le\frac{(u-v)^2}{v(1-v)}.
\end{equation}
This follows from $\log x\le x-1$ applied to the two likelihood ratios.
We next control how an arbitrary local observation history changes two
symmetric sign priors.

\begin{lemma}[Posterior contraction at every magnitude]\label{lem:posterior}
Fix $0<d\le\tfrac12$ and $0\le\eta<\tfrac12$. Under two experiments the sign priors are
$q_\pm=\tfrac12\pm\eta$, while the observation likelihoods conditional on the sign are the
same. Let $L\in[0,\infty]$ be the likelihood ratio of an arbitrary local history of one
coin under positive versus negative sign, let $\rho_\pm$ be its posterior sign
probabilities, and let $p_\pm=\tfrac12-d+2d\rho_\pm$ be the next predictive success
probabilities. If $0<L<\infty$, then
\begin{equation}\label{eq:posterior}
\rho_+-\rho_-=\frac{2\eta L}{L+(\tfrac14-\eta^2)(L-1)^2}\le2\eta .
\end{equation}
In every case,
\begin{equation}\label{eq:predictive}
\operatorname{kl}(p_+,p_-)\le\frac{32\,d^2\eta^2}{(1-4\eta^2)^2}.
\end{equation}
The same bound holds for $\operatorname{kl}(p_-,p_+)$. The constant $32$ is sharp: at $L=1$ the ratio of the two sides of
\eqref{eq:predictive} tends to one as $\eta\downarrow0$.
\end{lemma}

\begin{proof}
If $L\in\{0,\infty\}$, then $\rho_+=\rho_-\in\{0,1\}$, so $p_+=p_-$ and the divergence vanishes.
This case cannot occur for $d<\tfrac12$; for $d=\tfrac12$ it occurs after the first
observation. Now let $0<L<\infty$ and write $a=\tfrac12+\eta$, $b=\tfrac12-\eta$. Then
$a+b=1$ and $ab=\tfrac14-\eta^2=(1-4\eta^2)/4$. Bayes' formula gives
$\rho_+=aL/(b+aL)$ and $\rho_-=bL/(a+bL)$, so
\[
\rho_+-\rho_-=\frac{(a^2-b^2)L}{(b+aL)(a+bL)},\qquad
(b+aL)(a+bL)=ab(L-1)^2+(a+b)^2L .
\]
This proves \eqref{eq:posterior}. Also
\[
\rho_-(1-\rho_-)=\frac{abL}{(a+bL)^2},\qquad
\rho_+(1-\rho_+)=\frac{abL}{(b+aL)^2}.
\]
For fixed $v\in(0,1)$, the function $u\mapsto\operatorname{kl}(u,v)$ vanishes together with its
first derivative at $u=v$, and its second derivative is $1/\{u(1-u)\}$. Since
$p_-\le p_+$, Taylor's theorem gives
\[
\operatorname{kl}(p_+,p_-)\le\frac{(p_+-p_-)^2}{2\min_{\xi\in[p_-,p_+]}\xi(1-\xi)}.
\]
The same Taylor upper bound applies with the two arguments interchanged,
since its right side depends only on their squared difference and the
interval between them. Thus it controls both KL directions.
Every such $\xi$ equals $\tfrac12-d+2d\rho$ for some $\rho\in[\rho_-,\rho_+]\subset(0,1)$. Because
$d\le\tfrac12$,
\[
\xi(1-\xi)=\tfrac14-d^2(2\rho-1)^2\ge\tfrac14\{1-(2\rho-1)^2\}=\rho(1-\rho).
\]
By concavity, the minimum of $\rho(1-\rho)$ over $[\rho_-,\rho_+]$ is attained at an endpoint.
Since $p_+-p_-=2d(\rho_+-\rho_-)$, the preceding displays give
\[
\operatorname{kl}(p_+,p_-)\le
2d^2\max\Bigl\{\frac{(\rho_+-\rho_-)^2}{\rho_-(1-\rho_-)},
\frac{(\rho_+-\rho_-)^2}{\rho_+(1-\rho_+)}\Bigr\}
=\frac{8d^2\eta^2}{ab}\max\Bigl\{\frac{L}{(b+aL)^2},\frac{L}{(a+bL)^2}\Bigr\}.
\]
The maps $L\mapsto L/(b+aL)^2$ and $L\mapsto L/(a+bL)^2$ are maximized at $L=b/a$ and
$L=a/b$ respectively, both with value $1/(4ab)$. Therefore
$\operatorname{kl}(p_+,p_-)\le2d^2\eta^2/(ab)^2$, which is \eqref{eq:predictive}.

For sharpness, at $L=1$ we have $p_\pm=\tfrac12\pm2d\eta$ and
\[
\operatorname{kl}(p_+,p_-)=4d\eta\log\frac{1+4d\eta}{1-4d\eta}=32d^2\eta^2\{1+O(d^2\eta^2)\}.
\qedhere
\]
\end{proof}

Fix $h>0$ and take
\begin{equation}\label{eq:hard-pair}
 \eta_h(d)=\frac18\min\{1,h/d\},\qquad
 q_\pm(d)=\frac12\pm\eta_h(d).
\end{equation}
The two targets are centered at $1/2$ and separated by
\begin{equation}\label{eq:separation}
 \Delta_h=\theta_+-\theta_-=2\E_\mu\eta_h(D)=\frac14\om(h).
\end{equation}
Because $d\eta_h(d)\le h/8$, every next observation has predictive KL at
most $128h^2/225$ in either direction.

\begin{proposition}[Stopped-transcript bound]\label[proposition]{prop:transcript}
For every adaptive algorithm in any of the experiments of
\cref{thm:main}, its stopped transcript laws satisfy
\begin{equation}\label{eq:transcript}
 \KL(\mathsf P_+^{\cA},\mathsf P_-^{\cA})
       \le\frac{128}{225}h^2\E_+T.
\end{equation}
Consequently, if it estimates both targets to radius
$\varepsilon<\Delta_h/2$ with failure probability at most $\delta<1/2$,
\begin{equation}\label{eq:finite-lower}
 \E_+T\ge \frac{225}{128h^2}\kl(1-\delta,\delta).
\end{equation}
\end{proposition}
\begin{proof}
First reveal all magnitudes; an algorithm not using them is a special case.
Conditional on a realized transcript, action-selection and randomization
factors are common to the two hypotheses. Conditional on the hypothesis
and revealed magnitudes, the posterior of a coin's sign depends only on
its own observations: the likelihood factorizes by coin. Thus the local
posterior identity applies even after arbitrary revisiting and
interleaving. For a fresh coin the same assertion has $L=1$.
At $d=1/2$ the first observation has full support under both hypotheses,
and after it the sign is revealed, so later observations contribute zero KL.

A newly drawn magnitude has the same law in both experiments and contributes
zero KL. A selected magnitude is an algorithmic action and also contributes
zero conditional KL. The chain rule, first stopped at $T\wedge m$, bounds
transcript KL by $128h^2\E_+(T\wedge m)/225$. Padding the transcript after its
stopping time and taking increasing finite prefixes gives
\eqref{eq:transcript} by monotone convergence of expected work and the
finite-prefix characterization of relative entropy. If expected work is
infinite, the bound is vacuous. Interchanging the hypotheses gives the
reverse bound with $\E_-T$.

Thresholding $\widehat\theta$ at $1/2$ gives a test with both errors at
most $\delta$, since the estimation radius is strictly less than half
the separation. Data processing gives
$\KL\ge\kl(1-\delta,\delta)$, proving \eqref{eq:finite-lower}.
\end{proof}

The strict separation in this proposition is intentional. At exactly half
the separation, the constant output $1/2$ is within tolerance for both
members of the pair, so a test cannot be inferred without a limiting
argument over strictly separated pairs.

\section{A universal capped classifier}\label{sec:classifier}
This section does not assume a magnitude law. Fix a coin with
$P=1/2+(2S-1)d$, where $0<d\le1/2$ and $S\in\{0,1\}$.
For a dyadic cap $N=2^\ell$, examine prefixes of lengths $1,2,4,\ldots,N$.
At a prefix $n<N$, stop when
\begin{equation}\label{eq:threshold}
 \left|\overline X_n-\frac12\right|>t_{n,N},\qquad
 t_{n,N}=\sqrt{\frac{1+\log(N/n)}{n}},
\end{equation}
and output the side of $1/2$ occupied by the sample mean. If no early stop
occurs, output $Z_N=\ind\{\overline X_N>1/2\}$; a final tie is negative.
Let $\tau_N$ be the number of observations used.

Set
\begin{equation}\label{eq:m-definition}
 m_N(d)=\min\{N,d^{-2}[1+\logp(Nd^2)]\},\qquad
 \logp x=\max\{0,\log x\}.
\end{equation}

\begin{lemma}[Conditional error and work]\label[lemma]{lem:classifier}
For every dyadic $N\ge1$, both signs, and $0<d\le1/2$,
\begin{align}
 e_N(d,S):=\Prob(Z_N\ne S\mid d,S)
       &\le\min\{1,(Nd^2)^{-2}\},\label{eq:error}\\
 \E[\tau_N\mid d,S]&\le17m_N(d).\label{eq:work}
\end{align}
\end{lemma}
\begin{proof}
A wrong early stop at a fixed prefix $n$ requires a one-sided deviation
of at least $d+t_{n,N}$. Hoeffding's inequality and a union bound over
fixed prefixes, not a conditioning on the stopping event, give
\begin{align}
 e_N(d,S)
 &\le e^{-2Nd^2}+
       \sum_{\substack{n<N\\n\text{ dyadic}}}e^{-2n(d+t_{n,N})^2}
       \nonumber\\
 &\le e^{-2Nd^2}+e^{-2}
       \sum_{\substack{n<N\\n\text{ dyadic}}}(n/N)^2e^{-2nd^2}.
       \label{eq:error-sum}
\end{align}
For every dyadic sequence of positive $u$ values and every $v,b>0$,
\begin{equation}\label{eq:dyadic-integral}
 \sum_u u^v e^{-bu}
 \le \frac{\Gamma(v+1)}{(1-2^{-v})b^v}.
\end{equation}
Indeed, the disjoint intervals $[u/2,u]$ satisfy
$\int_{u/2}^u y^{v-1}e^{-by}\,\mathrm dy
\ge e^{-bu}u^v(1-2^{-v})/v$.
With $x=Nd^2\ge1$, the sum in \eqref{eq:error-sum} is bounded by
$x^{-2}(2/3)$, and $x^2e^{-2x}\le e^{-2}$. Hence the right side is at most
$(5/3)e^{-2}x^{-2}<x^{-2}$. The bound by one handles $x<1$.

The dyadic survival identity and the fact that a surviving prefix cannot
have crossed its threshold imply
\begin{align}
 \E\tau_N
 &=1+\sum_{n<N,\,n\text{ dyadic}}n\Prob(\tau_N>n)\nonumber\\
 &\le1+\sum_{n<N,\,n\text{ dyadic}}n
       \{\ind[d\le2t_{n,N}]+2e^{-nd^2/2}\}.\label{eq:survival}
\end{align}
For $x\ge1$, the inequality $d\le2t_{n,N}$ implies
$nd^2\le4(1+\log x)$: if $nd^2<1$ this is immediate, and otherwise
$\log(N/n)=\log(x/(nd^2))\le\log x$.
The corresponding dyadic sum of $n$ is at most
$8d^{-2}(1+\log x)$. Using \eqref{eq:dyadic-integral} with $v=1,b=1/2$
bounds the exponential part of \eqref{eq:survival} by $8d^{-2}$.
Since $1\le d^{-2}$, the total is at most
$17d^{-2}(1+\log x)$. Also $\tau_N\le N$ deterministically, and
$1+\log x\le x$ for $x\ge1$. For $x<1$, $m_N(d)=N$, completing the proof.
\end{proof}

\subsection{One geometric level law for every population}
On a common coin stream define
\begin{equation}\label{eq:deltas}
 \Delta_0=Z_1,\qquad \Delta_\ell=Z_{2^\ell}-Z_{2^{\ell-1}}
 \quad(\ell\ge1).
\end{equation}
A coupled difference is computed by running both classifiers on the same
prefix until both have stopped. Its work $W_\ell$ is the maximum of their
stopping times and hence at most their sum. Let
\begin{equation}\label{eq:p-level}
 r=2^{-1/2},\qquad p_\ell=(1-r)r^\ell,\qquad
 \Prob(J=\ell)=p_\ell.
\end{equation}
Choose $J$ independently of a fresh coin and its observations. For a finite
cutoff $L$, return
\begin{equation}\label{eq:Y-cut}
 Y_L=\frac{\Delta_J}{p_J}\ind\{J\le L\}.
\end{equation}
When $J>L$, no coin is sampled and the output and work are both zero.
Write $\mathcal W_L$ for actual work. For $L=\infty$, drop the indicator
and write $Y,\mathcal W$.

\begin{lemma}[Truncated inverse-margin control]\label[lemma]{lem:universal}
For every finite $L\ge0$, $N=2^L$, both signs, and $0<d\le1/2$,
\begin{align}
 \E[Y_L\mid d,S]&=\E[Z_N\mid d,S],\label{eq:tele}\\
 \E[Y_L^2\mid d,S]&\le85\min\{\sqrt N,d^{-1}\},\label{eq:Y2}\\
 \E[\mathcal W_L\mid d,S]&\le149\min\{\sqrt N,d^{-1}\}.
 \label{eq:Ywork}
\end{align}
Moreover,
\begin{equation}\label{eq:abs-bound}
 \sum_{\ell\ge0}\E[|\Delta_\ell|\mid d,S]
       \le C_0[1+\log(1/d)]
\end{equation}
for a universal $C_0$. Whenever $\E\log(1/D)<\infty$, $Y$ is integrable
and unbiased for $\theta$. This logarithmic condition is sufficient here;
no necessity claim for integrability alone is made. If $\E D^{-1}<\infty$,
its second moment is at most $85\E D^{-1}$ and its expected work at most
$149\E D^{-1}$.
\end{lemma}
\begin{proof}
The finite expectation telescopes. For $\ell\ge1$, a nonzero difference
requires at least one classifier to be wrong, so, with $N_\ell=2^\ell$,
\begin{equation}\label{eq:level-bounds}
 \E[\Delta_\ell^2\mid d,S]\le5\min\{1,(N_\ell d^2)^{-2}\},\qquad
 \E[W_\ell\mid d,S]\le34m_{N_\ell}(d).
\end{equation}
The same loose inequalities cover $\ell=0$. Consequently
\begin{align}
 \E[Y_L^2\mid d,S]
 &\le\frac5{1-r}\sum_{\ell=0}^L
       \sqrt{N_\ell}\min\{1,(N_\ell d^2)^{-2}\},\label{eq:Y2sum}\\
 \E[\mathcal W_L\mid d,S]
 &\le34(1-r)\sum_{\ell=0}^L
 N_\ell^{-1/2}m_{N_\ell}(d).\label{eq:Yworksum}
\end{align}
Split each sum at $N_\ell=d^{-2}$. Before that point, both summands are
$\sqrt{N_\ell}$ and their sum is bounded by
$(1-r)^{-1}\min\{\sqrt N,d^{-1}\}$.
After that point, the second-moment sum is a geometric tail with ratio
$2^{-3/2}$ and is at most $(1-r^3)^{-1}d^{-1}$.
Its full constant is bounded by
\[
 \frac5{1-r}\left(\frac1{1-r}+\frac1{1-r^3}\right)<85.
\]
For work, the first $N_\ell$ after the split has
$1<N_\ell d^2\le2$. Its tail is bounded by
\[
 d^{-1}\sum_{k\ge0}r^k(1+\log2+k\log2).
\]
Combining the prefix and tail gives the constant
\[
 34\left(2+\log2+\frac{r\log2}{1-r}\right)<149.
\]
If the cutoff precedes the split, only the smaller prefix bound is needed.
This proves \eqref{eq:Y2}--\eqref{eq:Ywork}.

Because $|\Delta_\ell|=\Delta_\ell^2$ for $\ell\ge1$, the number of
levels before the split is $O(1+\log(1/d))$, and the subsequent error
bounds form a convergent geometric series. This proves \eqref{eq:abs-bound}.
Absolute summability allows Fubini and telescoping to $S$, using
\eqref{eq:error}. Averaging over $(D,S)$ proves unbiasedness whenever the
stated logarithmic moment is finite. The inverse-moment condition implies
that condition. Finally take $L\to\infty$ in the nonnegative series for
second moment and work.
\end{proof}

This is the standard single-term randomized multilevel identity
\cite{vihola}; the common conditional bounds and their distributional
consequences are the key additional estimates. The estimator's output need
not lie in $[0,1]$. The interior-magnitude obstruction to nonnegative output is
given in \cref{sec:bounded}.

\section{Proof of the minimax law}\label{sec:minimax-proof}
\subsection{Exact calibration without a known margin law}
For $h=N^{-1/2}$ define
\begin{equation}\label{eq:bias-functional}
 b_\mu(h)=\E_\mu\min\{1,(h/D)^4\}\le\om(h).
\end{equation}
One call of the cutoff estimator returns $(Y_L,\mathcal W_L)$ with
\begin{equation}\label{eq:one-call}
 0\le\mathcal W_L\le N,\quad c:=\E\mathcal W_L\le149\Am(h),\quad
 v:=\Var(Y_L)\le85\Am(h),\quad |\E Y_L-\theta|\le b_\mu(h).
\end{equation}
Moreover $c>0$, because level zero has probability $1-r>0$ and costs one
observation. The following wrapper handles the unknown $c$ exactly.

\begin{lemma}[Two-success cost calibration]\label[lemma]{lem:calibration}
Suppose an independent call returns $(Z,C)$ with $0\le C\le N$,
$c=\E C>0$, $\E Z=\zeta$, and $\Var(Z)=v<\infty$.
Dependence of value and cost within a call is unrestricted.
Run independent pilot calls and mark each with conditional probability
$C_i/N$, using independent randomization. Stop at the second mark, at call
$K$. Discard all pilot values and average $M=K-1$ fresh independent
production values. If $\widehat\zeta$ is that average and $T$ is the total
pilot-plus-production cost, then
\begin{equation}\label{eq:calibration}
 \E\widehat\zeta=\zeta,\qquad
 \Var(\widehat\zeta)=\frac{vc}{N},\qquad
 \E T=4N-c\le4N.
\end{equation}
The procedure uses neither $c$ nor $v$.
\end{lemma}
\begin{proof}
The pilot marks are independent Bernoulli variables with parameter
$p=c/N>0$, so $K$ has negative-binomial probabilities
\[
 \Prob(K=k)=(k-1)p^2(1-p)^{k-2},\quad k\ge2,\qquad \E K=2/p.
\]
In particular $K<\infty$ almost surely, and the classical inverse-binomial
identity~\cite{haldane} is immediate:
\begin{equation}\label{eq:inverse-binomial}
 \E\frac1{K-1}=p^2\sum_{k=2}^\infty(1-p)^{k-2}=p.
\end{equation}
The event $\{K\ge i\}$ depends only on earlier calls and marks, so is
independent of $C_i$. Tonelli's theorem yields
\[
 \E\sum_{i=1}^K C_i=\sum_{i\ge1}\Prob(K\ge i)c=c\E K=2N.
\]
Production calls are independent of the pilot. Their mean total cost is
$c\E(K-1)=2N-c$, and conditional on $K$ the production average has mean
$\zeta$ and variance $v/(K-1)$. Total expectation and variance give
\eqref{eq:calibration}. For $p=1$, $K=2$ deterministically and the same
identities hold.
\end{proof}

\subsection{Distribution-blind upper bounds}
For $B\ge4$, choose the largest dyadic $N\le B/4$, set $L=\log_2N$,
and apply \cref{lem:calibration} to $(Z,C)=(Y_L,\mathcal W_L)$.
By \eqref{eq:one-call},
\begin{equation}\label{eq:blind-MSE}
 \E T\le4N\le B,\qquad
 \E(\widehat\theta-\theta)^2
 \le\om(h)^2+\frac{85\cdot149\Am(h)^2}{N}
 =12666\om(h)^2.
\end{equation}
Since $N>B/8$, \eqref{eq:scaling} gives the RMS upper bound, for example
with $C_R=\sqrt{8\cdot12666}$. For $1\le B<4$, output $1/2$ with no
observations; then $\om(B^{-1/2})=1$. This proves
\eqref{eq:blind-guarantee} without any knowledge of $\mu$ or $q$.

For confidence, choose the least odd integer
$k\ge(32/9)\log(1/\delta)$; then $k<5\log(1/\delta)$ for
$0<\delta\le1/4$. If $B/k\ge4$, use $k$ independent groups, each with
budget $B/k$, and return their median. A group's error exceeds
$\sqrt8 C_R\om(\sqrt{k/B})$ with probability at most $1/8$, by its
mean-square bound. Hoeffding's inequality bounds the probability that
at least half the groups fail by $e^{-9k/32}\le\delta$.
Costs add in expectation, and
\[
 \om(\sqrt{k/B})\le\sqrt5\,
 \om\!\left(\sqrt{\log(1/\delta)/B}\right).
\]
If $B/k<4$, output $1/2$; then
$\sqrt{\log(1/\delta)/B}>1/\sqrt{20}$, and
$\om(\sqrt{\log(1/\delta)/B})>1/\sqrt5$.
This handles the remaining case with a universal constant.

\paragraph{Implementation and scope.}
Pilot values are discarded; recycling them into the final average would
invalidate the independence calculation. Value and cost may be correlated
within any call. The random pilot and production counts are unbounded, so
the guarantee concerns expected observations, not a run-by-run work cap.
Projection of the final finite-budget estimate onto $[0,1]$ can only reduce
its error, but is not used in the variance identity.

\subsection{Matching lower bounds}
For RMS error use \eqref{eq:hard-pair} with $h=(4\sqrt B)^{-1}$.
By \cref{prop:transcript}, transcript KL is at most $8/225$, so Pinsker's
inequality gives total variation at most $2/15$.
Thresholding an estimator at $1/2$ gives a test whose error probabilities
sum to at most $8R^2/\Delta_h^2$, if its RMS error under both hypotheses
is at most $R$. Every test has error sum at least $1-2/15>1/2$.
Hence $R>\Delta_h/4=\om(h)/16\ge\om(B^{-1/2})/64$.
Taking the infimum proves the stated weak inequality.

For confidence write $K_\delta=\kl(1-\delta,\delta)$ and take
\[
 h=\frac{15}{16\sqrt2}\sqrt{K_\delta/B}.
\]
Then $225K_\delta/(128h^2)=4B>B$; by \cref{prop:transcript} no admissible
estimator has radius strictly smaller than $\Delta_h/2=\om(h)/8$.
On $0<\delta\le1/4$,
\begin{equation}\label{eq:confidence-kl}
 K_\delta=(1-2\delta)\log\frac{1-\delta}{\delta}
 \ge\frac14\log(1/\delta).
\end{equation}
Indeed $1-2\delta\ge1/2$ and $1-\delta\ge\sqrt\delta$.
Thus $h\ge15\sqrt{\log(1/\delta)/B}/(32\sqrt2)$.
Apply \eqref{eq:scaling} to obtain the coefficient $15/(256\sqrt2)$.
This completes the proof of \cref{thm:main}.

\subsection{Proof for a stochastic-dominance envelope}
The functions $d\mapsto\min\{1,h/d\}$,
$d\mapsto\min\{1,(h/d)^4\}$, and
$d\mapsto\min\{h^{-1},d^{-1}\}$ are decreasing. Their expectations for
any magnitude law in \eqref{eq:envelope-class} are therefore bounded by
their expectations under $\mu$. The distribution-blind upper algorithm is unchanged; its risk bound is
no larger than the envelope value. The lower bound follows by restricting to
$\cG(\mu)$, which is a subclass. The same reasoning with $d^{-1}$ and
\cref{thm:unbiased} proves \cref{cor:envelope}.

\begin{remark}[Why the statement uses error as a function of budget]
For highly irregular $\mu$, inverting $\om$ at two different constant
multiples of $\varepsilon$ need not preserve multiplicative constants in
work. Statements such as $B\asymp\om^{-1}(\varepsilon)^{-2}$ need a
regularity qualification or explicit constants inside the inverse.
The modulus bounds \eqref{eq:main-rms}--\eqref{eq:main-conf} do not have
this defect. Polynomial and logarithmic examples below can be inverted
with the stated asymptotic orders.
\end{remark}

\section{Proof of the unbiased characterization}\label{sec:unbiased-proof}
Sufficiency already follows from \cref{lem:universal}. The improved
$713I_\mu^2$ bound uses the modified classifier proved in \cref{sec:713}.
We prove the converse using only two-point inequalities and the stopped
KL chain rule; no differentiation of an estimator's expectation is needed.

\Needspace{10\baselineskip}
\begin{lemma}[Two-point variance inequality]\label{lem:hellinger-variance}
Let $\mathbb{P},\mathbb{Q}$ be probability measures on a common space, let
$H^2=1-\int\sqrt{d\mathbb{P}\,d\mathbb{Q}}$, and let $Y$ be a real statistic with finite second moment
under both. Then $\Delta=\mathbb{E}_{\mathbb{P}}Y-\mathbb{E}_{\mathbb{Q}}Y$ satisfies
\begin{equation}\label{eq:hellinger-variance}
\Delta^2(1-2H^2)\le4H^2\bigl(\operatorname{Var}_{\mathbb{P}}Y+\operatorname{Var}_{\mathbb{Q}}Y\bigr),
\end{equation}
Moreover, $2H^2\le\mathrm{KL}(\mathbb{P},\mathbb{Q})$, and the symmetric bound is
\begin{equation}\label{eq:symmetric-kl-hellinger}
8H^2\le\mathrm{KL}(\mathbb{P},\mathbb{Q})+\mathrm{KL}(\mathbb{Q},\mathbb{P}).
\end{equation}
\end{lemma}

\begin{proof}
Let $p,q$ be densities with respect to $\nu=\mathbb{P}+\mathbb{Q}$, and let
$c=(\mathbb{E}_{\mathbb{P}}Y+\mathbb{E}_{\mathbb{Q}}Y)/2$. Then
$\Delta=\int(Y-c)(\sqrt p-\sqrt q)(\sqrt p+\sqrt q)\,d\nu$. We use three facts:
$\int(\sqrt p-\sqrt q)^2\,d\nu=2H^2$; $(\sqrt p+\sqrt q)^2\le2(p+q)$; and
$\mathbb{E}_{\mathbb{P}}(Y-c)^2=\operatorname{Var}_{\mathbb{P}}Y+\Delta^2/4$, with the same identity under $\mathbb{Q}$.
Together with the Cauchy--Schwarz inequality they give
\[
\Delta^2\le\int(\sqrt p-\sqrt q)^2\,d\nu\cdot2\int(Y-c)^2(p+q)\,d\nu
=4H^2\Bigl(\operatorname{Var}_{\mathbb{P}}Y+\operatorname{Var}_{\mathbb{Q}}Y+\tfrac12\Delta^2\Bigr).
\]
Rearranging gives \eqref{eq:hellinger-variance}. Finally, by Jensen's inequality,
$-\tfrac12\mathrm{KL}(\mathbb{P},\mathbb{Q})=\mathbb{E}_{\mathbb{P}}\log\sqrt{q/p}\le\log(1-H^2)\le-H^2$.
For the symmetric inequality, use the pointwise bound
\[
(x-y)\log(x/y)\ge4(\sqrt{x}-\sqrt{y})^2,\qquad x,y>0.
\]
Indeed, putting $r=\sqrt{x/y}$, the difference is
$2y(r-1)(r+1)g(r)$, where
$g(r)=\log r-2(r-1)/(r+1)$ has $g(1)=0$ and
$g'(r)=(r-1)^2/[r(r+1)^2]\ge0$. Integrating the pointwise bound gives
\eqref{eq:symmetric-kl-hellinger}, with zero-density cases interpreted by limits.
\end{proof}


\begin{proof}[Converse in \cref{thm:unbiased}]
Suppose $Y$ is unbiased for every $q$, with finite uniform variance $V$ and
expected work $W$. Fix $h>0$ and $0<t\le h/8$, and consider
\[
 q_\pm(d)=\tfrac12\pm t\varphi_h(d),\qquad
 \varphi_h(d)=\min\{h^{-1},d^{-1}\}.
\]
Then $\eta(d)=t\varphi_h(d)\le t/h\le1/8$ and $d\eta(d)\le t$.
\Cref{lem:posterior} and the stopped-transcript argument give both
\begin{equation}\label{eq:converse-kl}
 \max\{\KL(\mathsf P_+^{\cA},\mathsf P_-^{\cA}),
        \KL(\mathsf P_-^{\cA},\mathsf P_+^{\cA})\}
 \le\kappa_t:=\frac{32t^2W}{(1-4t^2/h^2)^2}.
\end{equation}
Each direction uses its own expected work, bounded by $W$.
Unbiasedness gives mean separation $2t\Am(h)$.
By \eqref{eq:symmetric-kl-hellinger}, $H^2\le\kappa_t/4$.
Applying \cref{lem:hellinger-variance} with both variances at most $V$ yields
\[
 4t^2\Am(h)^2(1-\kappa_t/2)\le2\kappa_t V,
 \qquad
 \Am(h)^2(1-\kappa_t/2)
 \le\frac{16VW}{(1-4t^2/h^2)^2}.
\]
First let $t\downarrow0$, with $h$ fixed, giving $\Am(h)^2\le16VW$.
Then $\Am(h)\uparrow I_\mu$ as $h\downarrow0$ by monotone convergence.
This proves $I_\mu<\infty$ and $VW\ge I_\mu^2/16$.
The envelope class contains $\cG(\mu)$, so the same lower bound holds there.
\end{proof}

\paragraph{Sharpness of the universal coefficient.}
For $\mu=\delta_d$, $0<d\le1/2$, one fresh observation $X\in\{0,1\}$ gives
\begin{equation}\label{eq:fixed-gap-sharp}
 Y=\frac{X-1/2+d}{2d},\qquad \E Y=q(d),\qquad
 \sup_q\Var(Y)=\frac1{16d^2},\qquad W=1.
\end{equation}
Since $I_\mu=1/d$, equality holds. Hence $1/16$ cannot be increased in a
bound valid for every magnitude law. This does not assert equality for
every nondegenerate law, or that the parameter-free construction achieves
the optimal constant for each law.

\paragraph{Uniformity matters.}
The impossibility is over the whole sign class at a fixed $\mu$, or over
the stated envelope class. It does not say that each single population
with infinite inverse moment forbids every efficient unbiased estimator.
For a particular known population, the target might already be known.
Likewise, a small structured family of sign functions can be easier than
an arbitrary measurable $q$.

\subsection{Why a universal unbiased output must sometimes be negative}
\label{sec:bounded}
\begin{proposition}[No nonnegative unbiased output except at a deterministic endpoint]
\label[proposition]{prop:bounded}
If $\mu((0,1/2))>0$, no algorithm stopping almost surely for every $q$ and
always returning a nonnegative value is unbiased on $\cG(\mu)$.
For $\mu=\delta_{1/2}$ the obstruction disappears: one observation is an
unbiased $[0,1]$-valued estimator with unit work.
\end{proposition}
\begin{proof}
Under $q\equiv0$ the target is zero, so nonnegativity and unbiasedness force
the output to vanish almost surely. Compare this with
$q_1(d)=\ind\{d<1/2\}$, whose target is $\mu((0,1/2))>0$.
Interior coins give positive likelihood to every finite Bernoulli history
under both models; endpoint coins have the same deterministic sign zero.
Finite stopped transcript laws are equivalent. Intersecting with each
$\{T=n\}$ and taking a countable union gives equivalence of the full stopped
transcripts. Seeds and any revealed magnitudes have the same distribution.
Thus the output is also zero almost surely under $q_1$, a contradiction.
If $\mu=\delta_{1/2}$, the first observation is exactly the sign.
\end{proof}
This elementary support argument is not offered as a separate novelty
claim. It explains the signed outputs in the construction. Clipping is
allowed for a biased finite-accuracy estimator but invalidates exact
unbiasedness.

\section{A uniform variance--work upper constant of 713}\label{sec:713}
The finite-budget construction in \cref{sec:classifier,sec:minimax-proof}
uses the original threshold and constants $85,149$. For the unbiased
upper bound in \cref{thm:unbiased}, change the classifier threshold to
\[
 t_{n,N}=\sqrt{\frac{1/4+\log(N/n)}n}\qquad(n<N),
\]
keeping the same dyadic checkpoints, strict crossing rule, negative final
tie convention, shared-prefix increments, and independent level law
$p_\ell=(1-r)r^\ell$, $r=2^{-1/2}$. In this section $Z_N,\tau_N,
\Delta_\ell,Y_L,\mathcal W_L$ denote this modified construction only.
It observes no magnitude or sign and uses no distributional parameters.

\begin{theorem}[Separate conditional moment bounds]\label{thm:713}
For every $0<d\le1/2$, either sign, and finite $L$, $N=2^L$,
\[
 \E[Y_L^2\mid d,S]\le23\min\{\sqrt N,d^{-1}\},\qquad
 \E[\mathcal W_L\mid d,S]\le31\min\{\sqrt N,d^{-1}\}.
\]
If $I_\mu<\infty$, the untruncated estimator is square-integrable,
unbiased, and has finite expected actual work, with
\[
 \sup_q\E_{\mu,q}Y^2\le23I_\mu,\qquad
 \sup_q\E_{\mu,q}\mathcal W\le31I_\mu.
\]
Consequently the product of the separate variance and work suprema is
at most $713I_\mu^2$.
\end{theorem}

\paragraph{Error envelope.}
An incorrect early stop at prefix $n$ requires a one-sided deviation of
at least $d+t_{n,N}$. Hoeffding's inequality and a union bound over the
fixed dyadic prefixes yield, with $x=Nd^2$,
\[
 e_N\le e^{-2x}+e^{-1/2}N^{-2}
             \sum_{n<N,\ n\text{ dyadic}}n^2e^{-2nd^2}
 \le e^{-2x}+\tfrac23e^{-1/2}x^{-2}.
\]
For the second step integrate $u e^{-bu}$ over the disjoint intervals
$[n/2,n]$: each integral is at least $3n^2e^{-bn}/8$, so the sum is
at most $8/(3b^2)$. Since $x^2e^{-2x}\le e^{-2}$ for every $x>0$,
\[
 e_N\le\min\{1,A x^{-2}\},\qquad
 A=e^{-2}+\tfrac23e^{-1/2}<\tfrac17+\tfrac5{12}
 =\tfrac{47}{84}<\tfrac35.
\]
Here $e^2>7$ and $e^{1/2}>8/5$ follow from their positive power series.
For $\ell\ge1$, a nonzero increment requires a fine or coarse error;
the coarse scale has $x/2$. Hence
\begin{equation}\label{eq:713-increment}
 \E\Delta_\ell^2\le\min\{1,5A x^{-2}\}
             \le\min\{1,3x^{-2}\}.
\end{equation}
The same envelope covers $\ell=0$, because $d^2\le1/4$ and
$\Delta_0=Z_1\in\{0,1\}$.

\paragraph{Second moment.}
Put $d'=d/3^{1/4}$. Summing \eqref{eq:713-increment} divided by
$p_\ell$ and splitting at $2^\ell(d')^2=1$ gives
\[
 \E Y_L^2\le C_0\min\{\sqrt{2^L},1/d'\},\qquad
 C_0=\frac1{1-r}\left(\frac1{1-r}+\frac1{1-r^3}\right)<17.
\]
Before the split the terms are geometric with size $\sqrt{2^\ell}$;
after it their ratio is $r^3$. This includes cutoffs before the split.
The rational comparison $r<442/625$ and monotonicity of the displayed
expression prove $C_0<17$ by integer arithmetic. Since
$3^{1/4}<4/3$, the conditional constant is less than $68/3<23$.

\paragraph{Expected work.}
Take $c=13/7$. Survival at a dyadic prefix $n<N$ and $d>ct_{n,N}$
requires a one-sided deviation larger than $(1-1/c)d$. Therefore
\[
 \Prob(\tau_N>n)\le\ind\{d\le ct_{n,N}\}
            +\exp\{-2(1-1/c)^2nd^2\}.
\]
Suppose $x=Nd^2\ge1$ and write $u=nd^2$. The indicator implies
$u\le c^2[1/4+\log(x/u)]$. If $u<1$, then
$u\le1+c^2\log x$; if $u\ge1$, use $\log u\ge0$ and
$c^2/4=169/196<1$ for the same conclusion. Its dyadic contribution
is at most $2d^{-2}(1+c^2\log x)$.
For $b>0$, integration on $[n/2,n]$ gives
$\sum n e^{-bnd^2}\le2/(bd^2)$. The survival identity and
$1\le d^{-2}/4$ now give
\[
 \E\tau_N
 \le d^{-2}\left[\frac{125}{18}+\frac{338}{49}\log x\right]
 \le7d^{-2}(1+\log x).
\]
Together with $\tau_N\le N$, this proves $\E\tau_N\le7m_N(d)$.

On a shared stream the cap-$N$ threshold is at least the cap-$N/2$
threshold at every prefix before $N/2$. If the fine classifier stops
there, the coarse one has already stopped; at or after $N/2$ the coarse
one has reached its cap. Thus $\tau_{N/2}\le\tau_N$ pathwise and the
actual work of the coupled increment equals $\tau_N$. Level zero
costs one. It follows that
\[
 \E\mathcal W_L\le7(1-r)\sum_{\ell=0}^L
          2^{-\ell/2}m_{2^\ell}(d)
 \le7\left(2+\frac{\log2}{1-r}\right)
          \min\{\sqrt{2^L},d^{-1}\}.
\]
To obtain the last inequality, split at $2^\ell d^2=1$ as in
\cref{lem:universal}. The first tail value satisfies $1<2^\ell d^2\le2$;
the remaining sum is bounded by
$d^{-1}\sum_{k\ge0}r^k(1+\log2+k\log2)$.
Using $r<71/100$ and $\log2<7/10$ bounds the final constant by
$896/29<31$.

\paragraph{Untruncated limit and separate suprema.}
Integrate the conditional bounds over $(D,S)$. Monotone convergence of
the nonnegative moment and cost series gives finite second moment and
work when $I_\mu<\infty$. More explicitly,
\[
 \E[(Y-Y_L)^2]=\sum_{\ell>L}\frac{\E\Delta_\ell^2}{p_\ell}
       \longrightarrow0.
\]
Thus $Y_L\to Y$ in $L^2$ and $L^1$. Finite telescoping gives
$\E Y_L=\E Z_{2^L}$; the conditional error tends to zero for every
$d>0$, and bounded convergence proves $\E Y=\theta$.
The bounds hold for every $q$ with the same estimator and the same
$I_\mu$. Taking the two suprema separately and using
$\Var Y\le\E Y^2$ proves \cref{thm:713} and the improved upper
constant in \cref{thm:unbiased}. No interchange of a supremum with a
strict inequality is needed.

\paragraph{Formal scope.}
\leanref{Uniform713/FullBound.lean}{full\_bound\_713} proves the
untruncated estimator result in a conditional observation-kernel model,
including bounded independent observations with means $1/2\pm D$.
\leanref{ResearchRouteA/SeparateBounds.lean}{uniform\_family\_bounds}
supplies the separate-supremum conclusion for a fixed kernel and
estimator over a nonempty family of latent laws sharing $I_\mu$.
These results do not yet realize the finite-budget calibration wrapper
as a policy in the full fresh-coin experiment.


\section{Consequences for power and logarithmic margins}\label{sec:rates}
Let $F(t)=\mu((0,t])$. Stieltjes integration by parts gives, for $0<h<a$,
\begin{equation}\label{eq:cdf-formula}
 \om(h)=F(h)+h\int_{(h,a]}\frac1d\,\mu(\mathrm d d)
       =\frac ha+h\int_h^a\frac{F(t)}{t^2}\,\mathrm dt.
\end{equation}
The formula is valid with atoms. In particular,
\begin{equation}\label{eq:cdf-integrability}
 \Iinv=\frac1a+\int_0^a\frac{F(t)}{t^2}\,\mathrm dt.
\end{equation}
This integral, not a power exponent alone, determines exact unbiased
feasibility.

\subsection{The polynomial law as a corollary}
For $F(t)=(t/a)^\alpha$ on $(0,a]$ and $u=h/a<1$,
\begin{equation}\label{eq:power-modulus}
 \om(h)=
 \begin{cases}
  (\alpha u-u^\alpha)/(\alpha-1),&\alpha\ne1,\\
  u(1-\log u),&\alpha=1.
 \end{cases}
\end{equation}
Thus the main theorem immediately recovers the three regimes.
The usual polynomial-margin class allows a more general constant:
\begin{equation}\label{eq:power-class}
 \cG_{\alpha,C}=\{G\text{ on }[0,1]:
  \Prob_G(|P-1/2|\le t)\le Ct^\alpha,\quad0<t\le1/2\},
 \qquad C\ge2^\alpha.
\end{equation}
The last condition is necessary for nonemptiness.

\begin{corollary}[Sharp polynomial-margin complexity]\label[corollary]{cor:power}
Fix $\alpha>0$ and $C\ge2^\alpha$. For all sufficiently small
$\varepsilon>0$ and $0<\delta\le1/4$, the smallest uniform expected work
for an $(\varepsilon,\delta)$-valid estimator on $\cG_{\alpha,C}$ is
\begin{equation}\label{eq:power-complexity}
 \asymp_{\alpha,C}\log(1/\delta)
 \begin{cases}
 \varepsilon^{-2/\alpha},&0<\alpha<1,\\
 \varepsilon^{-2}\log^2(1/\varepsilon),&\alpha=1,\\
 \varepsilon^{-2},&\alpha>1.
 \end{cases}
\end{equation}
There is an unbiased estimator with uniformly finite variance and expected
work on this entire class if and only if $\alpha>1$. One common estimator
works for every such $\alpha,C$, without knowing either parameter.
\end{corollary}
\begin{proof}
Let $\mu_{\alpha,C}$ have CDF $\min\{Ct^\alpha,1\}$ on $(0,1/2]$.
It is supported on $(0,C^{-1/\alpha}]$, and
$\cG_{\alpha,C}=\cG_{\preceq\mu_{\alpha,C}}$ exactly.
Its modulus is \eqref{eq:power-modulus} with $a$ replaced by
$C^{-1/\alpha}$. Apply \cref{thm:main,cor:envelope} and invert the three
orders, with constants depending on $(\alpha,C)$.
The inverse moment is finite exactly when $\alpha>1$, so
\cref{thm:unbiased} gives both directions of the unbiased dichotomy.
Every polynomial envelope also has a finite uniform logarithmic moment;
thus the universal estimator is integrable and unbiased even for
$\alpha\le1$, when finite uniform variance and work cannot both hold.
\end{proof}

\subsection{Fixed-gap populations and rare-event scope}
\begin{corollary}[Gap classes]\label[corollary]{cor:gap}
For $0<\Delta\le1/2$, put
$\cG_\Delta=\{G:\Prob_G(|P-1/2|<\Delta)=0\}$.
For sufficiently small $\varepsilon$ and $0<\delta\le1/4$, its minimax
uniform expected work for $(\varepsilon,\delta)$ accuracy is
\[
 \asymp\frac{\log(1/\delta)}{\varepsilon^2\Delta^2},
\]
with universal comparison constants and a universal small-accuracy cutoff.
\end{corollary}
\begin{proof}
Here $\cG_\Delta=\cG_{\preceq\delta_\Delta}$ and
$\omega_{\delta_\Delta}(h)=\min\{1,h/\Delta\}$.
The envelope confidence law gives the result by direct inversion, taking
$\varepsilon$ smaller than its universal lower comparison constant.
\end{proof}
This is the balanced-population specialization of Lee--Valiant
\cite{leevaliant}. Their finer dependence on the positive fraction $\rho$
is important for rare events; it is outside our worst case over all $q$.
Our hard pairs keep every sign prior within $1/8$ of $1/2$, so both targets
are near $1/2$. The localized theory is developed in \cref{sec:rare}; a sharp
characterization for every magnitude law remains open.

\subsection{New boundary regimes invisible to a single exponent}
For $\beta\ge0$, define an actual CDF on $(0,a]$ by
\begin{equation}\label{eq:log-cdf}
 F_\beta(t)=\frac{t/a}{[1+\log(a/t)]^\beta}.
\end{equation}
It is increasing, vanishes at zero, and equals one at $a$.
With $L_h=1+\log(a/h)$, its modulus is exactly
\begin{equation}\label{eq:log-modulus}
 \omega_\beta(h)=\frac ha
 \begin{cases}
  1+\dfrac{L_h^{1-\beta}-1}{1-\beta},&\beta\ne1,\\[6pt]
  1+\log L_h,&\beta=1,
 \end{cases}\qquad 0<h<a.
\end{equation}
This follows from \eqref{eq:cdf-formula} by the substitution
$u=1+\log(a/t)$.

\begin{corollary}[Logarithmically refined critical law]\label[corollary]{cor:log}
For the fixed magnitude law \eqref{eq:log-cdf}, or its stochastic-dominance
envelope class, the optimal expected work at small accuracy
$\varepsilon$ is, up to constants depending on $\beta,a$,
\begin{equation}\label{eq:log-rates}
 \log(1/\delta)\times
 \begin{cases}
 \varepsilon^{-2}[\log(1/\varepsilon)]^{2(1-\beta)},&0\le\beta<1,\\
 \varepsilon^{-2}[\log\log(1/\varepsilon)]^2,&\beta=1,\\
 \varepsilon^{-2},&\beta>1.
 \end{cases}
\end{equation}
Uniformly finite variance and work for unbiased estimation are possible
exactly when $\beta>1$, in which case
$\E D^{-1}=\beta/[a(\beta-1)]$.
\end{corollary}
\begin{proof}
Equation \eqref{eq:log-modulus} and \cref{thm:main,cor:envelope} give the
error law. Its monotone inverse has the displayed orders: if
$\omega(h)\asymp h[\log(1/h)]^\gamma$ with fixed $\gamma\ge0$, then
$h\asymp\varepsilon/[\log(1/\varepsilon)]^\gamma$; replacing the log by
$\log\log$ gives the boundary case. Finally,
\[
 \E D^{-1}=\frac1a\left(1+\int_1^\infty u^{-\beta}\,\mathrm du\right),
\]
which is finite exactly when $\beta>1$. Apply \cref{thm:unbiased}.
\end{proof}

For $\beta>1$, this class permits efficient unbiased estimation even
though it satisfies no bound $F(t)\le C t^\alpha$ with any $\alpha>1$ and
finite $C$. Indeed, $t^{1-\alpha}/[1+\log(a/t)]^\beta\to\infty$.
The earlier exponent boundary is therefore a corollary about entire
polynomial classes, not the general structural criterion.

\begin{table}[tb]
\centering\small
\begin{tabularx}{\linewidth}{@{}lXl@{}}\toprule
Magnitude law near zero & Optimal work at fixed confidence & Efficient unbiased estimator\\\midrule
$F(t)\asymp t^\alpha$, $\alpha<1$ & $\varepsilon^{-2/\alpha}$ & No\\
$F(t)\asymp t$ & $\varepsilon^{-2}\log^2(1/\varepsilon)$ & No\\
$F(t)\asymp t^\alpha$, $\alpha>1$ & $\varepsilon^{-2}$ & Yes\\
$F(t)\asymp t/[\log(1/t)]^\beta$, $\beta<1$ &
$\varepsilon^{-2}[\log(1/\varepsilon)]^{2(1-\beta)}$ & No\\
$F(t)\asymp t/\log(1/t)$ & $\varepsilon^{-2}[\log\log(1/\varepsilon)]^2$ & No\\
$F(t)\asymp t/[\log(1/t)]^\beta$, $\beta>1$ & $\varepsilon^{-2}$ & Yes\\\bottomrule
\end{tabularx}
\caption{Consequences of the modulus law. ``Efficient unbiased'' means both
variance and expected work bounded uniformly over the corresponding sign or
envelope class, not a claim for every individual population. For high
confidence, multiply work by $\log(1/\delta)$.}
\label{tab:rates}
\end{table}

\paragraph{An iterated-log example.}
The same calculation extends further. Set
$L_1(t)=1+\log(a/t)$, $L_2(t)=1+\log L_1(t)$, and
$F(t)=(t/a)/(L_1(t)L_2(t)^\gamma)$ with $\gamma\ge0$.
Then
\[
 \omega(h)=\frac ha\left(1+\int_1^{L_2(h)}v^{-\gamma}\,\mathrm dv\right).
\]
At $\gamma=1$, the fixed-confidence work is
$\asymp\varepsilon^{-2}[\log\log\log(1/\varepsilon)]^2$, for sufficiently
small $\varepsilon$. Unbiased feasibility again changes at $\gamma>1$.
This is an immediate instance of the general integral criterion, not an
additional assumption or an independently designed algorithm.

\section{Persistence under fixed finite-order density smoothness}
\label{sec:smooth}
The arbitrary-sign model makes no smoothness assumption on $G$. We now
show that the critical logarithmic penalty is not solely caused by a
jump in the density at the threshold.

For an integer $s\ge1$, write
$\|f\|_{C^s}=\max_{0\le k\le s}\|f^{(k)}\|_\infty$.

\begin{theorem}[Smooth-density robustness]\label{thm:smooth}
For each fixed integer $s\ge1$, there exists a finite $M_s$ such that,
for every fixed $M\ge M_s$, the class of $C^\infty$ probability densities
supported in $[1/4,3/4]$ and satisfying $\|f\|_{C^s}\le M$ has minimax
expected complexity
\[
 \asymp_{s,M}\varepsilon^{-2}\log^2(1/\varepsilon)\log(1/\delta)
\]
for sufficiently small $\varepsilon$ and $0<\delta\le1/4$.
No unbiased estimator has both uniformly finite variance and uniformly
finite expected work on this class.
\end{theorem}
\begin{proof}
Fix the local cutoff $a_0=1/4$. Choose a fixed even $C^\infty$ density $g$ supported strictly inside
$(-a_0,a_0)$, with $g\ge c_g>0$ on $[-a_0/2,a_0/2]$. The centered coordinate is
$x=p-1/2$. For sufficiently small $h>0$, set
\begin{equation}\label{eq:smooth-pair}
 t_h=h^{1/(s+1)},\qquad
 f_\pm(1/2+x)=g(x)\left(1\pm\frac{2hx}{x^2+t_h^2}\right).
\end{equation}
The perturbation is odd, so both densities integrate to one. Its absolute
size is at most $h/t_h=h^{s/(s+1)}$, which is less than $1/4$ for small
$h$, so the densities are nonnegative. The common magnitude density is
$2g(d)$ and the conditional sign perturbation is
\[
 \eta_h(d)=\frac{hd}{d^2+t_h^2},\qquad
 \eta_h(d)\le\frac{h}{2t_h}\le\frac18,\qquad d\eta_h(d)\le h.
\]
The posterior lemma therefore gives predictive KL at most
$(8192/225)h^2$ per observation, uniformly over histories and selected coins.
The target separation satisfies
\begin{align*}
 \theta_+-\theta_-
 &=4h\int_0^{a_0} g(d)\frac{d}{d^2+t_h^2}\,\mathrm dd\\
 &\ge2c_g h\log\frac{(a_0/2)^2+t_h^2}{t_h^2}
 \ge c_s h\log(1/h)
\end{align*}
for all sufficiently small $h$, with $c_s>0$; for a fixed bump it can be chosen proportional to
$(s+1)^{-1}$. This is not uniform as $s\to\infty$.

For $R_t(x)=x/(x^2+t^2)=t^{-1}R_1(x/t)$,
\[
 \|hR_{t_h}^{(k)}\|_\infty
 \le C_k h t_h^{-(k+1)}
 =C_k h^{(s-k)/(s+1)}\le C_k,\qquad 0\le k\le s.
\]
Leibniz's rule bounds every $\|f_\pm\|_{C^s}$ by a fixed $M_s$ independent
of $h$. They are individually $C^\infty$ because $t_h>0$ and $g$ is a
smooth compactly supported function. Uniform bounds on derivatives of
all orders at once are not asserted.

Choose $h$ of order $\varepsilon/\log(1/\varepsilon)$ with a sufficiently
large fixed factor so that the separation exceeds $2\varepsilon$.
The stopped-transcript testing argument proves the lower bound. For the
upper bound, $\|f\|_\infty\le M$ implies
$\Prob(|P-1/2|\le t)\le2M t$, so \cref{cor:power} applies.
The unbiased impossibility follows by averaging a putative uniform
finite-variance, finite-work estimator and contradicting the lower bound
at fixed confidence as $\varepsilon\downarrow0$.
\end{proof}

\paragraph{Which smoothness is being controlled.}
This theorem concerns the density of the latent $P$. It is not a statement
about smoothness of a regression function $y\mapsto\E[X\mid Y=y]$ in
observable covariates. Smooth cross-state structure can enable information
sharing, as in kernel methods~\cite{chen}. Neither this theorem nor the
next one rules out such improvements.

The construction does not settle analogous bounds for fixed analytic or
Gevrey regularity classes; those require separate assumptions and arguments.

\section{An observed-outer-state embedding}\label{sec:outer}
A latent coin and an observable outer state are different experiments.
The following construction supplies a lower bound in the latter without
silently identifying the two.

Let $Y$ be uniform on $[0,1]$. The algorithm may observe independent outer
states and may also select any state at which to request a fresh inner
observation. Let $X\in\{-1,1\}$ and put
$\theta=\Prob(\E[X\mid Y]>0)$. There is no regularity assumption connecting
the conditional means at different outer states.

\begin{proposition}[Finite-cell observable-state lower bound]
\label[proposition]{prop:outer}
Fix a probability law $\mu$ on $(0,1/2]$, $0<h\le1/2$, and
$0<\delta\le1/8$. There is a finite partition of $[0,1]$ and a family of
deterministic, cellwise-constant conditional Bernoulli models such that:
(a) their magnitude laws obey the envelope $F_\mu$;
(b) their construction has a scale separation
$\Delta\ge\omega_\mu(h)/8$; and (c) every estimator valid to radius
$\Delta/4$ and failure probability $\delta$ uniformly over this family
has worst-case expected work at least
\begin{equation}\label{eq:outer-lower}
 \frac{225}{128h^2}\kl(1-2\delta,2\delta).
\end{equation}
This lower bound permits arbitrary selected-state queries and revisiting.
\end{proposition}
\begin{proof}
Round each $d\le h$ up to $h$. Round larger $d$ upward to the next value
in the grid $h,2h,4h,\ldots$, capped at $a$. This gives a finitely supported
law $\mu_h$ on $\{d_1,\ldots,d_J\}$, with weights $w_j$, which is
stochastically no smaller than $\mu$. Above $h$ the rounding increases a
magnitude by at most a factor two; below $h$ the clipped integrand stays
one. Hence
\[
 \omega_{\mu_h}(h)\ge\tfrac12\omega_\mu(h).
\]
For each $j$, create $M$ cells of lengths $w_j/M$ and assign unknown fixed
signs $s_{jk}\in\{0,1\}$. On cell $(j,k)$ let
\[
 B\mid Y\sim\Ber\bigl(1/2+(2s_{jk}-1)d_j\bigr),\qquad X=2B-1.
\]
Every deterministic sign array has magnitude law $\mu_h$, meets the
envelope, and has target
$\theta(s)=\sum_j(w_j/M)\sum_k s_{jk}$.

For the testing proof only, draw signs independently under two product
priors with probabilities $1/2\pm\eta_h(d_j)$ from
\eqref{eq:hard-pair}. Their target means are separated by
\[
 \Delta=\tfrac14\omega_{\mu_h}(h)\ge\tfrac18\omega_\mu(h).
\]
The sum of squared target weights is
$\sum_j w_j^2/M\le1/M$. Weighted Hoeffding therefore gives
\[
 \Prob_\pm(|\theta(s)-\E_\pm\theta(s)|>\Delta/8)
       \le2\exp(-M\Delta^2/32).
\]
Choose an integer $M\ge32\Delta^{-2}\log(2/\delta)$.
Uniform estimation error at most $\Delta/4$, together with the displayed
concentration event, makes thresholding at $1/2$ a test with each error
at most $2\delta$.

The signs are independent under either prior, so the local posterior
lemma applies separately to every cell, regardless of the adaptively
selected query points. Each inner observation has predictive KL at most
$128h^2/225$. A uniform expected-work bound for deterministic sign arrays
also bounds the prior-averaged expected work. Apply the KL chain rule
and binary data processing to obtain \eqref{eq:outer-lower}.
\end{proof}

The partition size depends on the desired scale and confidence. This is
a minimax construction in a class allowing arbitrarily fine finite
partitions, not an asymptotic lower bound for one fixed finite list of
cells. Since the outer distribution itself stays uniform on $[0,1]$, it
establishes the relevant obstruction even with observable covariates,
provided no cross-state regularity is imposed. The original upper
algorithm can ignore the labels and use fresh outer states, so the
envelope modulus remains the matching error scale for this
unstructured class.

\Needspace{13\baselineskip}
\subsection{A continuous-magnitude subclass satisfying the nested assumptions}
\begin{proposition}[Continuous-magnitude observed-state obstruction]\label{prop:continuous}
Fix the local cutoff $a_0=1/4$. Consider observed outer states $U\sim\mathrm{Unif}[0,1]$ and binary inner observations $X\in\{-1,1\}$ with
\[
|\E[X\mid U]|=2D(U),\qquad D(U)\sim\mathrm{Unif}(0,a_0].
\]
No cross-state regularity is assumed. There is a class of deterministic conditional models of this form whose minimax expected observation cost for error $\varepsilon$ with fixed failure probability $1/16$ is
\[
\asymp\varepsilon^{-2}\log^2(1/\varepsilon).
\]
The lower bound allows arbitrary selected-state queries and revisiting. All models satisfy the bounded-density and normalized conditional-moment assumptions 2.1 and 2.4 of Giles--Haji-Ali \cite{gileshajiali}, with constants independent of the error scale.
\end{proposition}
\begin{proof}
Fix $0<h<a_0$ and partition $(0,a_0]$ into the first band $(0,h]$ and consecutive dyadic bands with right endpoints $2h,4h,\ldots$, clipping the final endpoint at $a_0$. Let band $j$ have endpoints $\ell_j,b_j$ and mass $w_j=(b_j-\ell_j)/a_0$. Partition $[0,1]$ into $M$ equal-weight cells per band, each of length $w_j/M$. Within every cell in band $j$, let $D(U)$ increase linearly across $(\ell_j,b_j]$. Define its values at the finitely many endpoints to be positive and in the assigned band. Consequently the magnitude law is exactly uniform on $(0,a_0]$, not a discretization. The partition and the map $D$ may be disclosed to the algorithm.

Assign each cell an unknown deterministic sign $s_{jk}\in\{0,1\}$. On that cell let $X=2B-1$ for $B\sim\Ber(1/2+(2s_{jk}-1)D(U))$. The target is
\[
\theta(s)=\sum_j\frac{w_j}{M}\sum_{k=1}^M s_{jk}.
\]
For the proof only, choose independent cell signs under two priors with probabilities $1/2\pm\eta_j$, where $\eta_j=\frac18\min\{1,h/b_j\}$. Their target means have separation
\[
\Delta_h=2\sum_jw_j\eta_j\ge\frac18\omega_\mu(h),\qquad \mu=\mathrm{Unif}(0,a_0].
\]
Indeed, within the first band the clipped integrand is one; within every other band $b_j\le2D(U)$. Every queried magnitude in band $j$ satisfies $D(U)\eta_j\le h/8$.

Repeated queries in a cell may use different magnitudes. Nevertheless, its sign has a single common likelihood ratio under the two priors, obtained by multiplying its queried sign-conditional likelihoods. The posterior identity and the predictive bound therefore still apply. Each observation contributes at most $128h^2/225$ KL, even under arbitrary query selection.

Weighted Hoeffding gives
\[
\Prob_\pm\bigl(|\theta(s)-\E_\pm\theta(s)|>\Delta_h/8\bigr)
\le2e^{-M\Delta_h^2/32}.
\]
For $M\ge32\Delta_h^{-2}\log(2/\delta)$, an estimator with uniform radius $\Delta_h/4$ and failure probability $\delta\le1/8$ gives a test of these priors with each error at most $2\delta$. The stopped-transcript argument gives expected work at least
\[
\frac{225}{128h^2}\kl(1-2\delta,2\delta).
\]
For the uniform magnitude law,
$\omega_\mu(h)=(h/a_0)[1+\log(a_0/h)]$. Set $h$ of order $\varepsilon/\log(1/\varepsilon)$ with a sufficiently large fixed constant to obtain the lower rate. The upper rate follows by ignoring state labels and using the preceding threshold estimator. The class permits arbitrarily many cells; it is not a fixed finite-partition asymptotic claim.

Finally, $\Var(X\mid U)=1-4D(U)^2\ge3/4$. Their normalized magnitude is
\[
Z=\frac{|\E[X\mid U]|}{\sqrt{\Var(X\mid U)}}
=\frac{2D}{\sqrt{1-4D^2}}.
\]
It has density $[2a_0(1+z^2)^{3/2}]^{-1}$ on its support, bounded by $1/(2a_0)$. Moreover, $|X-\E[X\mid U]|\le3/2$, so for every finite $q>2$ its standardized conditional $q$th absolute moment is bounded by $3^{q/2}$. These verify the two stated assumptions uniformly.
A uniform RMS guarantee $\varepsilon$ gives error at most
$4\varepsilon$ with failure probability at most $1/16$, by Markov's
inequality. Applying the fixed-confidence lower bound at $4\varepsilon$
yields the same RMS rate. The RMS upper bound follows from
\cref{thm:main}. This also proves \cref{cor:nested}.
\end{proof}

\paragraph{Interpretation.}
This identifies a hard subclass satisfying the explicit density and moment assumptions behind the classical adaptive nested-simulation upper rate. It does not rule out improvements from smooth cross-state regression structure, a known parametric model, or restrictions to rare-event populations. A restriction on the density of the latent magnitude is different from a restriction on the function of observed covariates.
